\chapter{FD2NN 4f Multi-Layer Sweep}

\section{Purpose}

The spatial D2NN branch had shown meaningful detector-side gains. The corrected
4f FD2NN question was therefore straightforward: if the same Raw-RP objective is
applied to a five-layer Fourier-plane stack, can the architecture reproduce the
same scale of detector gain?

\section{Sweep Setup}

The \texttt{0405-fd2nn-4f-sweep-pitchrescale} family varies $f \in \{10, 25\}$\,mm and
$z \in \{1, 3, 5\}$\,mm while keeping the loss fixed to
\texttt{focal\_raw\_received\_power}.

\runfigure{0405-fd2nn-4f-sweep-pitchrescale/fd2nn_sweep_figures/fig1_pib_bar.png}
  {FD2NN PIB bar chart}
  {Detector-side focal PIB at 10\,$\mu$m across the six 4f FD2NN settings.}

\runfigure{0405-fd2nn-4f-sweep-pitchrescale/fd2nn_sweep_figures/fig2_heatmap.png}
  {FD2NN heatmap}
  {Heatmap of the six tested $(f,z)$ combinations.}

\runfigure{0405-fd2nn-4f-sweep-pitchrescale/fd2nn_sweep_figures/fig3_loss_curves.png}
  {FD2NN loss curves}
  {Training curves for the six 4f FD2NN configurations.}

\runfigure{0405-fd2nn-4f-sweep-pitchrescale/fd2nn_sweep_figures/fig4_tp_co.png}
  {FD2NN throughput and CO}
  {Throughput and output-plane CO across the 4f sweep.}

\runfigure{0405-fd2nn-4f-sweep-pitchrescale/fd2nn_sweep_figures/fig5_regime.png}
  {FD2NN regime figure}
  {Interpretive diagram for the $f$--$z$ trade-off.}

\runfigure{0405-fd2nn-4f-sweep-pitchrescale/fd2nn_sweep_figures/fig6_phase_masks.png}
  {FD2NN masks}
  {Learned phase masks for the six sweep points.}

\begin{table}[H]
  \centering
  \begin{tabular}{lccccc}
    \toprule
    \smalltablehead{Config} & \smalltablehead{$\Delta x_f$ ($\mu$m)} & \smalltablehead{$\PIB@10\,\mu$m} & \smalltablehead{$\Delta$ vs base} & \smalltablehead{$\TP$} & \smalltablehead{Detector interpretation} \\
    \midrule
    baseline & -- & 0.5752 & -- & -- & reference \\
    f10mm\_z1mm & 7.57 & 0.3774 & -0.1978 & 1.0000 & below baseline \\
    f10mm\_z3mm & 7.57 & 0.2902 & -0.2850 & 1.0000 & below baseline \\
    f10mm\_z5mm & 7.57 & 0.2510 & -0.3242 & 1.0000 & below baseline \\
    f25mm\_z1mm & 18.92 & 0.5856 & +0.0104 & 1.0000 & near-tie / slight gain \\
    f25mm\_z3mm & 18.92 & 0.4302 & -0.1450 & 1.0000 & below baseline \\
    f25mm\_z5mm & 18.92 & 0.4192 & -0.1560 & 1.0000 & below baseline \\
    \bottomrule
  \end{tabular}
  \caption{Sweep-level focal-bucket summary for the 4f FD2NN study. Because
  throughput is essentially unity in every tested configuration, the detector
  interpretation follows the same ordering as the focal-bucket column: only
  \texttt{f25mm\_z1mm} slightly exceeds the baseline.}
\end{table}

\section{The $f$--$z$ Dilemma}

The sweep reveals a structural trade-off. Smaller focal lengths encourage
stronger diffraction mixing but make the Fourier plane harder to sample and more
sensitive to defocus. Larger focal lengths are physically safer, but the
mixing becomes weak enough that the detector gain nearly disappears. The tested
best point, \texttt{f25mm\_z1mm}, should therefore be read as a near-tie rather than a
clean success.

\begin{keyinsight}
The 4f FD2NN branch does not collapse because of energy loss; throughput is
essentially unity everywhere. It underperforms because the Fourier-plane
architecture does not deliver enough useful detector-side shaping under the
tested geometry.
\end{keyinsight}

This sets up the final interpretive chapter of the technical core: what kind of
fundamental limit does the 4f architecture appear to face?
